% =========================================================
% Bridge: From Real Analysis to Metric Spaces
% Source: volume-iii/analysis/metric-spaces/notes/metric-spaces/
%
% Purpose: Make explicit which results from real analysis depend
%          on |·| as a metric vs. on properties specific to ℝ.
% =========================================================

\subsubsection{Porting Results from $\mathbb{R}$ to Metric Spaces}
\label{sec:real-analysis-bridge}

\begin{tcolorbox}[colback=gray!6, colframe=gray!40, arc=2pt,
  left=8pt, right=8pt, top=6pt, bottom=6pt,
  title={\small\textbf{Why abstract this?}},
  fonttitle=\small\bfseries]
Real analysis proved many theorems about sequences and functions in
$\mathbb{R}$ using $|x - y|$ to measure closeness.
When you look back at those proofs carefully, you find two very different
kinds of steps:

\medskip
\noindent
\textbf{Steps that used only \emph{distance}:}
``choose $N$ so that for all $n \ge N$, $|x_n - x| < \varepsilon$'' ---
this is just asking that $d(x_n, x) < \varepsilon$.

\medskip
\noindent
\textbf{Steps that used properties of $\mathbb{R}$ specifically:}
``since $(x_n)$ is bounded and monotone, it converges'' --- this
requires the Monotone Convergence Theorem, which uses the order structure
of $\mathbb{R}$ and fails in $\mathbb{R}^2$.

\medskip
\noindent
The theorems in the first category are exactly those that port
to arbitrary metric spaces.
The ones in the second category require extra structure (completeness, compactness, etc.)
and either fail in general or need a new proof strategy.
\end{tcolorbox}

\bigskip

\begin{center}
\small
\textbf{What Ports Directly vs.\ What Needs Extra Structure}
\smallskip

\begin{tabular}{p{5.2cm} p{4.4cm} p{4.4cm}}
\toprule
\textbf{Result from $\mathbb{R}$} & \textbf{Metric space version} & \textbf{What is needed} \\
\midrule
$x_n \to x$ means $|x_n - x| \to 0$ &
$x_n \to x$ means $d(x_n, x) \to 0$ &
Only the metric \\[6pt]

Convergent sequences are Cauchy &
Same: convergent $\Rightarrow$ Cauchy &
Only triangle inequality \\[6pt]

Cauchy sequences converge &
\emph{Fails in general} (e.g.\ $\mathbb{Q}$) &
Completeness \\[6pt]

Bounded monotone sequences converge &
No analogue &
Order structure of $\mathbb{R}$; not portable \\[6pt]

Convergent sequences are bounded &
Same: convergent $\Rightarrow$ bounded &
Only triangle inequality \\[6pt]

Bolzano--Weierstrass: bounded sequences have convergent subsequences &
\emph{Fails in general} ($\ell^2$ has bounded non-compact sets) &
Compactness \\[6pt]

Continuous functions on $[a,b]$ attain their extrema &
\emph{Fails in general} &
Compactness of domain \\[6pt]

Closed and bounded $\Leftrightarrow$ compact &
Only closed $\Leftarrow$ compact in general &
Finite dimension ($\mathbb{R}^n$) \\
\bottomrule
\end{tabular}
\end{center}

\begin{remark*}[How to read the proofs in this chapter]
Whenever you see a proof that uses only the metric axioms (M1)--(M4)
and no other structure of the space, that proof works in \emph{every}
metric space automatically.
Whenever you see a proof that assumes completeness or compactness,
pause and ask: is this telling me something about $\mathbb{R}$ specifically,
or about a class of metric spaces?
This habit of mind is what abstract analysis is training you to develop.
\end{remark*}

\begin{example}[A proof that ports perfectly]
\label{ex:convergent-implies-cauchy-port}
In $\mathbb{R}$ the proof that convergent implies Cauchy went:

\medskip
Given $\varepsilon > 0$, choose $N$ so that $n \ge N \Rightarrow |x_n - L| < \varepsilon/2$.
Then for $m, n \ge N$: $|x_m - x_n| \le |x_m - L| + |L - x_n| < \varepsilon/2 + \varepsilon/2 = \varepsilon$.

\medskip
\noindent
Every step here uses only the triangle inequality for $|\cdot|$,
which is precisely the triangle inequality (M4) for the metric $d(x,y) = |x-y|$.
The proof in a general metric space is \emph{word for word} the same with $d$ in
place of $|\cdot|$:

\medskip
Given $\varepsilon > 0$, choose $N$ so that $n \ge N \Rightarrow d(x_n, x) < \varepsilon/2$.
Then for $m, n \ge N$: $d(x_m, x_n) \le d(x_m, x) + d(x, x_n) < \varepsilon$.

\medskip
\noindent
This is why metric spaces exist: to identify exactly which results
are consequences of \emph{having a distance} rather than of
\emph{being the real line}.
\end{example}
